\documentclass[10pt,a4paper]{amsart}
\usepackage[margin=19mm]{geometry}
\usepackage{amsmath,amssymb,mathtools}
\usepackage[scr=rsfs,cal=euler]{mathalfa}
\usepackage{xfrac}
\usepackage[unicode,hidelinks,bookmarks=false]{hyperref}
\newcommand{\ie}{i.e.\ }
\newcommand{\cf}{cf.\ }
\newcommand{\resp}{respectively\ }
\newcommand{\Subsection}{\subsection}
\providecommand{\nobreakdash}{\nobreak\hbox{-}}
\setlength{\emergencystretch}{2em}
\multlinegap0pt
\usepackage{bm}
\usepackage{bbm}
\usepackage{dsfont}
\usepackage{xspace}
\usepackage{cancel}
\usepackage{mathtools}
\usepackage{amsthm}
\makeatletter
\@ifundefined{lemma}{\newtheorem{lemma}{Lemma}}{}
\@ifundefined{theorem}{\newtheorem{theorem}{Theorem}}{}
\makeatother
\newcommand{\norm}[1]{\|#1\|}
\title[State-based approach to the numerical solution of Dirichlet ]{State-based approach to the numerical solution of Dirichlet boundary optimal control problems for the Laplace equation}
\author{Ulrich Langer, Richard L\"oscher, Olaf Steinbach, Huidong Yang}
\date{Source: arXiv:2507.11646v1 (CC BY 4.0)}
\begin{document}
\maketitle

\noindent\emph{Source and licence.} State-based approach to the numerical solution of Dirichlet boundary optimal control problems for the Laplace equation. Version arXiv:2507.11646v1 (CC BY 4.0), distributed under the Creative Commons Attribution 4.0 International licence, as recorded in the arXiv licence metadata for this exact version and shown on the abstract page https://arxiv.org/abs/2507.11646v1. Full rights notice: licenses/arxiv-2507.11646-CC-BY-4.0.txt.\par\smallskip

\noindent\textit{Editorial scope.} Counted unit: the Theorem whose source label is \texttt{Eqn:Error Y} in the article (source0.tex lines 1202--1217, source label \texttt{Eqn:Error Y}), delivered together with the article's own complete original proof of it (source0.tex lines 1219--1271, one \texttt{proof} environment of 53 lines). No mathematics is written, repaired or re-derived: statement and proof are the article's own text line for line. Required context retained uncounted: Eqn:Regularization error L2 (lines 302--305); Eqn:Regularization error Y (lines 314--319); Eqn:VF 12 (lines 924--931); Eqn:FEM (lines 946--953); Eqn:Cea (lines 976--990); Eqn:Error p (lines 1026--1035). Every definition, display and label of those lines is retained unchanged. Omitted from this package and not redistributed: 1--301, 306--313, 320--923, 932--945, 954--975, 991--1025, 1036--1201, 1218--1218, 1272--1911. Nothing inside a retained excerpt is omitted or altered: every retained body is delivered exactly as the article's lines are written, so no omission receipt of any kind accompanies this package. Citations: the retained excerpts carry no citation command at all, so no bibliography entry of the article is transcribed and no reference is redistributed. Reference adaptation: none was needed and none was made; every reference inside the delivered unit resolves inside it, so no unresolved reference or citation is produced. Printed slips kept literal: no printed word, symbol, hypothesis or number of the counted statement or of its proof was altered. Layout and class: the article's own class is \texttt{article} with packages including a4, paralist, graphicx, psfrag, pgfplotstable, subcaption, amsmath, amsthm, amsfonts, booktabs, diagbox, multirow, hyperref, changes; the wrapper is the lane's pinned \texttt{amsart} editorial wrapper at 10pt on A4, which restates the theorem-like environments the retained text needs and adds the article's own macro definitions that the retained text needs (\texttt{\textbackslash norm}), with extra wrapper packages (bm, bbm, dsfont, xspace, cancel, mathtools). The delivered numbering is the unit's own. Assets: no retained span contains a figure environment, a graphics-inclusion command, a PDF-page inclusion, a table or a TikZ picture; no image file is copied into this package, the package declares no asset dependency and no image file is redistributed with it. Licence: version arXiv:2507.11646v1 of this article is distributed under the Creative Commons Attribution 4.0 International licence, as recorded in the arXiv licence metadata for this exact version and shown on the abstract page https://arxiv.org/abs/2507.11646v1. A full rights notice is delivered at licenses/arxiv-2507.11646-CC-BY-4.0.txt. Characters: every retained excerpt is pure ASCII, so the delivered engine, XeLaTeX with its default Latin Modern text font, reports no missing character.
\begin{equation}\label{Eqn:Regularization error L2}
  \| y_\varrho - \overline{y} \|_{L^2(\Omega)} \leq
  \| \overline{y} \|_{L^2(\Omega)}, 
\end{equation}

\begin{equation}\label{Eqn:Regularization error Y}
  \| y_\varrho - \overline{y} \|_{L^2(\Omega)} \leq
  \varrho^{1/2} \, \| \nabla \overline{y} \|_{L^2(\Omega)}, \quad
  \| \nabla (y_\varrho - \overline{y}) \|_{L^2(\Omega)} \leq
  \| \nabla \overline{y} \|_{L^2(\Omega)}, 
\end{equation}

\begin{equation}\label{Eqn:VF 12}
  \begin{array}{lclcl}
    a(y_\varrho,y) + b(y,p_\varrho)
    & = & \langle \overline{y},y \rangle_{L^2(\Omega)}
    && \forall y \in V, \\[1mm]
    b(y_\varrho,q) & = & 0 & & \forall q \in X,
  \end{array}
\end{equation}

\begin{equation}\label{Eqn:FEM}
  \begin{array}{lclcl}
    a(y_{\varrho h},y_h) + b(y_h,p_{\varrho h})
    & = & \langle \overline{y},y_h \rangle_{L^2(\Omega)}
    && \forall y_h \in V_h, \\[1mm]
    b(y_{\varrho h},q_h) & = & 0 & & \forall q_h \in X_h.
  \end{array}
\end{equation}

\begin{lemma}
  Let $(y_\varrho,p_\varrho) \in V \times X$ and
  $(y_{\varrho h},p_{\varrho h}) \in V_h \times X_h$ be the unique
  solutions of the variational formulations \eqref{Eqn:VF 12} and
  \eqref{Eqn:FEM}, respectively. Then 
  %there holds Cea's lemma, i.e.,
  the error estimate
  \begin{equation}\label{Eqn:Cea}
    \| y_\varrho -y_{\varrho h}\|_V
    \, \leq \, 2 \,
    \| y_\varrho-y_h\|_V  +
    \frac{1}{\sqrt{\varrho}} \, \| p_\varrho-q_h \|_X
  \end{equation}
  holds for all $(y_h,q_h) \in Y_h \times X_h$.
\end{lemma}

\begin{lemma}
  Assume that $\Omega$ is either a convex Lipschitz domain, or a bounded
  domain with smooth boundary.
  Let $(y_\varrho,p_\varrho) \in V \times X$ be the unique solution of the
  variational formulation \eqref{Eqn:VF 12}. Then $p_\varrho\in H^2(\Omega)$ and
  \begin{equation}\label{Eqn:Error p}
    \inf_{q_h\in X_h}\| p_\varrho - q_h \|_X \leq \norm{p_\varrho-I_h p_\varrho}_X\leq 
    c \, h \, \| y_\varrho - \overline{y} \|_{L^2(\Omega)} .
  \end{equation}
\end{lemma}

\begin{theorem}
  Let $y_{\varrho h} \in V_h$ be the unique solution of the variational
  formulation \eqref{Eqn:FEM}. Assume that $\Omega$ is either a convex Lipschitz domain, or a bounded domain with smooth boundary. For $\overline{y} \in Y$ there holds the error estimate
  \begin{equation}\label{Eqn:Error Y}
    \| y_{\varrho h} - \overline{y} \|_{L^2(\Omega)} \leq
    c \, \sqrt{h^2 + \varrho} \, \| \nabla \overline{y} \|_{L^2(\Omega)} =
    \widetilde{c} \, h \, \| \nabla \overline{y} \|_{L^2(\Omega)},
  \end{equation}
  when choosing $\varrho = h^2$, while for $\overline{y} \in L^2(\Omega)$
  we have
   \begin{equation}\label{Eqn:Error L2}
    \| y_{\varrho h} - \overline{y} \|_{L^2(\Omega)} \leq
    c \, \sqrt{1 + h^2 \varrho^{-1}} \, \| \overline{y} \|_{L^2(\Omega)} =
    \widetilde{c} \, \| \overline{y} \|_{L^2(\Omega)}.
  \end{equation}
\end{theorem}

\begin{proof}
  For $\overline{y} \in Y$ and using \eqref{Eqn:Error p} and
  \eqref{Eqn:Regularization error Y} we first have
  \[
    \| p_\varrho - I_h p_\varrho \|_X \leq
    c \, h \, \| y_\varrho - \overline{y} \|_{L^2(\Omega)} \leq
    c \, h \, \varrho^{1/2} \, \| \nabla \overline{y} \|_{L^2(\Omega)} .
  \]
  For $y_h^* \in V_h$ being the finite element approximation of
  $y_\varrho \in Y$ we now have, using standard finite element error
  estimates and \eqref{Eqn:Regularization error Y},
  \begin{eqnarray*}
    \| y_\varrho - y_h^* \|_V^2
    & = & \| y_\varrho - y_h^* \|_{L^2(\Omega)}^2 + \varrho \,
          \| \nabla (y_\varrho - y_h^*) \|^2_{L^2(\Omega)} \\
    & \leq & c \, \Big[ h^2 + \varrho \Big] \,
             \| \nabla y_\varrho \|^2_{L^2(\Omega)} \, \leq \,
             c \, \Big[ h^2 + \varrho \Big] \,
             \| \nabla \overline{y} \|^2_{L^2(\Omega)},
  \end{eqnarray*}
  Cea's lemma \eqref{Eqn:Cea} now gives
  \[
    \| y_\varrho - y_{\varrho h} \|^2_{L^2(\Omega)}
    \leq 4 \, \| y_\varrho - y_h^* \|_V^2 + \frac{2}{\varrho} \,
    \| p_\varrho - I_h p_\varrho \|_X^2 \leq
    c \, \Big[ h^2 + \varrho \Big] \,
    \| \nabla \overline{y} \|^2_{L^2(\Omega)} .
  \]
  Hence, and using \eqref{Eqn:Regularization error Y} we
  obtain \eqref{Eqn:Error Y},
  \[
    \| y_{\varrho h} - \overline{y} \|_{L^2(\Omega)}^2
    \leq 2 \, \| y_{\varrho h} - y_\varrho \|^2_{L(\Omega)} +
    2 \, \| y_\varrho - \overline{y} \|^2_{L^2(\Omega)} \\
    \leq c \, \Big[ h^2 + \varrho \Big] \,
    \| \nabla \overline{y} \|^2_{L^2(\Omega)} .
  \]
  For $\overline{y} \in L^2(\Omega)$ we conclude,
  now using \eqref{Eqn:Regularization error L2},
   \[
    \| p_\varrho - I_h p_\varrho \|_X \leq
    c \, h \, \| y_\varrho - \overline{y} \|_{L^2(\Omega)} \leq
    c \, h \, \| \overline{y} \|_{L^2(\Omega)} ,
  \]
  and
  \[
    \| y_\varrho - y_h^* \|_V^2 \leq c \, \Big[ h^2 + \varrho \Big] \,
    \| \nabla y_\varrho \|^2_{L^2(\Omega)}
    \leq c \, \Big[ h^2 + \varrho \Big] \, \varrho^{-1} \,
    \| \overline{y} \|^2_{L^2(\Omega)} .
  \]
  With this, \eqref{Eqn:Error L2} follows.
\end{proof}

\end{document}
